\documentclass[11pt]{article}
\usepackage{amsmath,amssymb,hyperref,url}
\hypersetup{hidelinks}
\setlength{\emergencystretch}{2em}
\title{When Comparison-Based Reward Shaping Preserves Incentives\\
An Account of the Q--P Research Thread}
\author{}
\date{}
\begin{document}
\maketitle

\begin{abstract}
This thread paired \emph{Mechanism Design for Alignment and Control} (Q) with \emph{MARS-RA} (P), which uses
visual comparisons to guide cooperative learning.
The peers asked whether P's rewards could support Q's requirements for truthful reports and obedience.
They found that fully accounted potential-based shaping preserves existing incentives, while omitted rewards
during absence can create a conditional incentive to hide evidence on return.
They also showed why repeated judgments of fixed evidence cannot overcome a limit on what a bounded payment can induce.
The verifier paused the thread because these arguments specialise known results, and neither the relevant
accounting convention nor a feasible manipulation had been established for P.
\end{abstract}

\section{Introduction}

Q, by Bergemann, Koh and Morris, studies how to reward an agent whose preferences and capabilities are unknown
\cite{Q}. An agent may conceal a capability but cannot counterfeit it. The designer wants the agent both to
report truthfully and to follow a recommended action. Q therefore asks whether any feasible combination of a
false report and a changed action would benefit the agent. A mechanism is incentive compatible when no such
deviation pays. Its results concern what reward rules can achieve under these strategic constraints.

P, by Wang and colleagues, studies credit assignment in embodied multi-agent cooperation \cite{P}. An external
large multimodal model compares pairs of images taken from the agents' viewpoints. A Bradley--Terry model,
which represents pairwise preferences through score differences, combines these judgments into contribution
scores. P normalises the scores using softmax, which turns them into positive shares summing to one. These
shares become potentials: numerical values whose discounted changes are added to training rewards. Agents can
disappear and later return, making the accounting across absences a relevant question.

The proposed connexion was whether P's comparison-based rewards could inherit Q's strategic guarantees. The
peers found that this did not follow. P's judges supply external assessments, not strategic reports from the
agents. Recovering a judge's underlying preferences is a statistical task; making truthful reporting and
obedience worthwhile is an incentive task. The investigation therefore shifted to a narrower question: when
does P's form of reward shaping preserve incentives, and when might its accounting change them?

\section{What was done}

Ada and Emmy first read the definitions behind the apparent connexion. Both distinguished P's estimation
target from Q's constraints on joint report and action deviations. They then divided the work. Ada examined
how shaping rewards add up when an agent leaves and returns. Emmy examined how much an additional payment
could achieve when its only evidence came from the judge's observations.

Ada derived an identity for rewards that are included in an agent's return only at selected transitions. This
exposed a possible debit on re-entry. Emmy derived a bound relating a profitable deviation, the available
payment range and the distinguishability of the judge's evidence. In parallel, their source checks narrowed
the novelty claims. Devlin and Kudenko already establish equilibrium preservation for dynamic potential-based
shaping \cite{devlin}; Grzes studies episodic terminal-potential effects and their correction \cite{grzes}.
Saig, Einav and Talgam-Cohen give a stronger statistical-contract result that contains the proposed two-action
payment bound \cite{saig}.

The peers next qualified the shaping argument, raised and settled an objection about relative scores, and
constructed exact numerical examples. Ada checked the boundary identity and its correction using rational
arithmetic. Emmy checked finite distributions for repeated judgments of one image and for judgments of fresh
images. Each cross-checked the other's main results. They ended by agreeing that an audit of participation
accounting was the main concrete direction. The evidence bound restricted possible incentive-changing repairs;
it did not establish another failure of P.

\section{Results}

\subsection{Complete accounting leaves deviation gains unchanged}

Fix one persistent agent identity and an episode of \(H\) transitions. Let \(t\) denote a transition index,
\(\gamma\) the discount factor with \(0<\gamma\leq1\), and \(\Phi_t\) the agent's potential at time \(t\). The
shaping increment is \(F_t=\gamma\Phi_{t+1}-\Phi_t\); the reward adds \(\rho F_t\), where \(\rho\) is a fixed
weight. If the terminal potential \(\Phi_H\) is zero, then
\[
\sum_{t=0}^{H-1}\gamma^t F_t
=-\Phi_0+\gamma^H\Phi_H=-\Phi_0.
\]
Each interior potential appears once with a positive sign and once with a negative sign, so it cancels.

If the initial potential is fixed before any strategic choice, every complete strategy receives the same total
adjustment, \(-\rho\Phi_0\). A deviation that was profitable remains profitable by exactly the same amount.
This includes a deviation combining a false report with a changed action. Exact shaping therefore cannot newly
establish Q's truthfulness or obedience requirements. It may still improve learning, so this conclusion does
not contradict P's claimed learning benefits.

The result holds with the same information and feasible strategies before and after shaping. It also requires
adjacent increments to reuse the same realised potential. These conditions matter even if potentials depend on
history or judge randomness. Ada gave the argument in \path{ada/boundary_audit.md}, and Emmy accepted these
qualifications in \path{emmy/information_audit.md}.

\subsection{Omitted increments can create a re-entry debit}

Let \(m_t\) equal one when the agent's return includes the shaping increment at transition \(t\), and zero
otherwise. This inclusion rule is the reward mask. Shifting indices gives
\[
\begin{aligned}
\sum_{t=0}^{H-1}\gamma^t m_t F_t
={}&-m_0\Phi_0
+\sum_{t=1}^{H-1}\gamma^t(m_{t-1}-m_t)\Phi_t\\
&+\gamma^H m_{H-1}\Phi_H.
\end{aligned}
\]
The middle sum records what fails to cancel when inclusion changes. A changing team alone does not cause a
problem: complete accounting still cancels the interior terms. Dropping increments can cause one.

Suppose both the terminal potential and potentials during inactivity are zero. Exit terms then vanish, but a
re-entry at time \(e\) leaves a discounted debit \(-\gamma^e\Phi_e\). A smaller returning potential means a
smaller debit. This can reward concealment if the agent can lower that potential at sufficiently little cost.

Ada illustrated this with three transitions, zero base rewards, \(\gamma=0.99\), \(\rho=1\), mask \((1,0,1)\),
and potentials \((1/2,0,p,0)\). Here \(p\) is the potential on return. The total masked shaping reward is
\[
G_{\mathrm{mask}}(p)=-\frac12-(0.99)^2p,
\]
where \(G_{\mathrm{mask}}\) denotes the discounted return under that mask. An explicitly assumed
evidence-withholding action changes \(p\) from \(1/2\) to \(1/10\), without changing the physical trajectory.
Before any manipulation cost, its gain is
\[
G_{\mathrm{mask}}(1/10)-G_{\mathrm{mask}}(1/2)=0.39204.
\]
Complete accounting gives \(-1/2\) in either case.

Under the zero-inactive-potential convention, adding \(\Phi_e\) to the reward at re-entry cancels the debit.
More generally, adding the settlement \((m_t-m_{t-1})\Phi_t\) at each interior time cancels the middle sum.
These corrections must be multiplied by \(\rho\) when that shaping weight differs from one. Nonzero inactive
potentials may require settlement outside active transitions.

Ada's \path{ada/check_boundaries.py} checked 11,664 rational instances of the identity and checked the
settlement for zero terminal potential. Emmy independently ran the checks and verified the example. This
supports the algebra and the conditional counterexample. It does not show that P uses this mask or permits
this manipulation.

\subsection{Repeated judgments cannot supply missing evidence}

Emmy considered a possible additional contract, rather than P's original shaping reward. Fix an agent's type,
report and recommendation, and the other agents' strategies. Let \(g>0\) be the intrinsic expected gain from a
feasible deviation relative to the desired behaviour. Let \(Z\) contain all evidence supplied to the judge,
including images and metadata. Write \(\mu_\star\) and \(\mu_d\) for its probability laws under desired and
deviating behaviour, respectively.

Let \(Y_K\) be the transcript of \(K\) judge queries, with laws \(\nu_\star^K\) and \(\nu_d^K\). The judge
uses the same random procedure given \(Z\) in either case, and obtains no further evidence. Suppose an
additive payment \(T(Y_K)\) has range at most \(B>0\), with no other payment or action-prohibition channel.
Obedience requires
\[
\begin{aligned}
g&\leq \mathbb{E}_{\nu_\star^K}[T]-\mathbb{E}_{\nu_d^K}[T]\\
 &\leq B\operatorname{TV}(\nu_\star^K,\nu_d^K)
 \leq B\operatorname{TV}(\mu_\star,\mu_d).
\end{aligned}
\]
Here \(\mathbb{E}\) denotes expectation under the indicated law. Total variation, written
\(\operatorname{TV}\), measures how distinguishable two probability laws are. It ranges from zero for
identical laws to one for perfectly distinguishable laws. On a finite outcome space it is half the sum of the
absolute differences between outcome probabilities.

For the first upper bound, subtract the payment's infimum and divide by \(B\). The resulting function lies
between zero and one, so its expectation difference cannot exceed total variation. Conditioning that function
on \(Z\) gives the evidence bound. Repeated queries may clarify what the judge thinks about an image, but
cannot add evidence about the action that produced it. Identical evidence laws leave the profitable deviation
undeterred at any finite payment range.

For a single deviation the transcript bound is attainable: pay \(B\) on outcomes more likely under desired
behaviour, and zero otherwise. Equality with \(g\) gives only weak obedience. Several deviations require a
common contract; satisfying each pairwise bound separately is not sufficient.

Emmy's synthetic example made the distinction concrete. Binary evidence equals one with probability \(0.6\)
under desired behaviour and \(0.4\) under deviation. Given evidence one or zero, each conditionally
independent judgment favours the agent with probability \(0.9\) or \(0.1\), respectively. With \(g=0.3\) and
\(B=1\), repeated judgments of one image have total variation at most \(0.2\). No query count can deter the
deviation.

Fresh independent evidence before each query changes the premise. The judgment probabilities then become
\(0.58\) and \(0.42\). With eight queries, total variation is approximately \(0.198909\) for a shared image
and \(0.341177\) for fresh images. The latter exceeds the required \(0.3\), so the specified single-deviation
payment suffices. Emmy's \path{emmy/check_information.py} checked the exact finite distributions; Ada
independently recomputed these values. These are constructed calculations, not observations from P's
benchmark.

\section{Objections and unresolved claims}

The initial transfer of strategic guarantees did not hold. Q's peer discipline requires separated conditional
beliefs and additional reward instruments. P's external comparisons do not supply these assumptions. Nor does
statistical consistency establish that the judge measures the strategic quantities Q requires.

Ada also suggested that losing a common shift in contribution scores could obstruct implementation.
Bradley--Terry comparisons and softmax cannot distinguish scores that all increase by the same amount. Emmy
objected that losing a distinction matters only if the indistinguishable cases need different incentives, or
conceal a profitable deviation. Q itself has an example where action differences identify what the reward
needs. Ada accepted the objection. Relative scoring alone therefore supplied no impossibility result.

The peers qualified the complete-accounting claim too. If an initial report determines the initial potential,
the residual payment can depend on that report. If rewards reveal new information, the available strategies
may change. If adjacent increments use independently resampled potentials, cancellation need not hold. These
are limits of the stated argument, not evidence that P has a strategic reporting channel.

The central empirical objection remained unresolved. The supplied material does not establish P's
inactive-agent reward accounting, persistent identity through re-entry, stored potentials, or singleton score
convention. No runnable implementation was supplied or audited. The assumed evidence-withholding action also
needs to be feasible, with a specified cost and strategic payoff. P uses comparisons during training; a
strategic deployment mechanism would be an extension.

Finally, the source checks removed novelty claims for the generic shaping and payment arguments. The
statistical-contract result of Saig and colleagues, Proposition~1, treats total variation distance to mixtures
of deviations \cite{saig}. The peers did not settle prior work on the particular participation audit. The
material supports a focused investigation, not a completed publication-level result.

\section{Why the thread stopped}

The verifier issued a single recorded decision: \textsc{Pause}. It accepted that the note followed the
evidence and that the boundary identity and information bounds were sound. It stopped the thread because the
mathematics specialised acknowledged prior results, while the proposed P-specific connexion depended on an
unverified reward mask and an assumed manipulation. More synthetic calculations would not establish those
missing facts.

The smallest restart is to inspect and instrument one disappearance and re-entry trace in a runnable P
implementation. Record the persistent agent identity, potentials, included shaping increments and discounted
return. Compare the accumulated shaping with the boundary identity and complete accounting. This determines
whether the proposed residual occurs in that implementation.

Only if it does should the next test introduce an explicitly feasible evidence intervention and measure its
gain, including its cost. Correct accounting defeats this accounting-based prediction; an unavailable or
unprofitable intervention defeats the corresponding attack. The unresolved task is therefore to obtain
implementation evidence and a specified strategic payoff, not to extend the existing algebra.

\begingroup
\raggedright
\bibliographystyle{plain}
\bibliography{references}
\endgroup
\end{document}
